\begin{abstract}
Probabilistic coherence spaces are the standard denotational model of
higher-order probabilistic programs. We ask which of their elements are
mixtures of deterministic ones. For each type we take the mixtures of its
certain states, the elements with values in $\{0,1\}$, as its probability
calculus. At every formula of multiplicative--additive linear logic over
finite data types, the certain states are cliques of Girard's coherence
space and the calculus lies inside the probabilistic coherence space. The
two agree at first order and differ from second order on. For a program
with two functional arguments, the calculus is the stable set polytope of
a graph whose clique-constrained polytope is the coherence space, and five
web points separate them, with the bounds of the KCBS inequality. For the
dual of a tensor of two additive types, the two coincide iff one of the
two coherence graphs has no induced four-cycle, and the gap transfers to
every formula in which such a tensor occurs negatively. Linear negation
exchanges the calculus with a local calculus cut out by the certain states
of the dual, and the coherence space is the self-dual set between them.
The calculus is not compositional: first-order types are tight, but their
tensor is not. At the dual of the second-order type, the gap is Bell's
theorem, witnessed by the Popescu--Rohrlich box. All results are
mechanized in Lean~4.

\keywords{Probabilistic coherence spaces \and Linear logic \and
Probabilistic programs \and Perfect graphs \and Contextuality \and Lean}
\end{abstract}
